% Ad Atticum 14.6 (ad-atticum-14-06)
% Source: https://raw.githubusercontent.com/PerseusDL/canonical-latinLit/master/data/phi0474/phi057/phi0474.phi057.perseus-lat1.xml
% Fetched by scripts/fetch_latin.py

% Dateline (Perseus): Scr. Fundis prid. Id. Apr. a. 710 (44).

\ciceroLetterOpener{CICERO ATTICO S. D.}

\ciceroSection{1} pridie Idus Fundis accepi tuas litteras cenans. primum igitur melius esse, deinde meliora te nuntiare. odiosa illa enim fuerant, legiones venire. nam de Octavio susque deque. exspecto quid de Mario; quem quidem ego sublatum rebar a Caesare. Antonio conloquium cum heroibus nostris pro re nata non incommodum. sed tamen adhuc me nihil delectat praeter Idus Martias. nam quoniam Fundis sum cum Ligure nostro, discrucior Sextili fundum a verberone Curtilio possideri.

\ciceroSection{2} quod cum dico, de toto genere dico. quid enim miserius quam ea nos tueri propter quae illum oderamus? etiamne consules et tribunos pl. in biennium quos ille voluit? nullo modo reperio quem ad modum possim πολιτεύεσθαι . nihil enim tam σόλοικον quam tyrannoctonos in caelo esse, tyranni facta defendi. sed vides consules, vides reliquos magistratus, si isti magistratus, vides languorem bonorum. exsultant laetitia in municipiis. dici enim non potest quanto opere gaudeant, ut ad me concurrant, ut audire cupiant mea verba de re p. nec ulla interea decreta. sic enim πεπολιτεύμεθα ut victos metueremus. haec ad te scripsi apposita secunda mensa; plura et πολιτικώτερα postea, et tu quid agas quidque agatur.
